\documentclass[10pt,a4paper]{amsart}
\usepackage[margin=19mm]{geometry}
\usepackage{amsmath,amssymb,mathtools}
\usepackage[scr=rsfs,cal=euler]{mathalfa}
\usepackage{xfrac}
\usepackage[unicode,hidelinks,bookmarks=false]{hyperref}
\newcommand{\ie}{i.e.\ }
\newcommand{\cf}{cf.\ }
\newcommand{\resp}{respectively\ }
\newcommand{\Subsection}{\subsection}
\providecommand{\nobreakdash}{\nobreak\hbox{-}}
\setlength{\emergencystretch}{2em}
\multlinegap0pt
\usepackage{thmtools}
\usepackage{thm-restate}
\DeclareMathOperator*{\E}{\mathbb{E}}
\newcommand{\Prob}{\mathbb P}
\DeclareMathOperator{\TV}{d_{\mathrm{TV}}}
\theoremstyle{plain}
\newtheorem{theorem}{Theorem}
\newtheorem{proposition}[theorem]{Proposition}
\newtheorem{lemma}[theorem]{Lemma}
\newtheorem{corollary}[theorem]{Corollary}
\theoremstyle{definition}
\newtheorem{definition}[theorem]{Definition}
\newtheorem{assumption}[theorem]{Assumption}
\theoremstyle{remark}
\newtheorem{remark}[theorem]{Remark}
\title[Query Lower Bounds for Diffusion Sampling]{Query Lower Bounds for Diffusion Sampling: Proposition E.1}
\author{Zhiyang Xun, Eric Price}
\date{Source: arXiv:2604.10857v2 (CC BY 4.0)}
\begin{document}
\maketitle

\noindent\emph{Source and licence.} The delivered work is the article shown in the title above, version arXiv:2604.10857v2 (CC BY 4.0), distributed under the Creative Commons Attribution 4.0 International licence, as recorded in the arXiv licence metadata for this exact version and shown on the abstract page saved with this item. The full licence notice, which carries the licence URL and the address of that abstract page, is the bundled file \texttt{licenses/arxiv-2604.10857-CC-BY-4.0.txt}.\par\smallskip

\noindent\textit{Editorial scope.} Counted unit: the article's Proposition E.1 (source0.tex lines 1703--1720, source label \texttt{prop:bounded-reweighting}), printed by the article in its appendix E; statement and proof are the article's own text line for line, in the article's own notation ($\nu_S$, $\nu^\eta_S$, $\mathcal D^\eta$, $n_{\max}$, $C_{\mathrm{wt}}$, $B_d$), with no line omitted, substituted or re-flowed. Required context retained uncounted: the article's hard-construction setup (source0.tex lines 655--701), its dense-codebook lemma (lines 1377--1386) and its level-sensitive null-coupling proposition (lines 1538--1550). Editorial formatting only: an AMS \texttt{amsart} preamble that restates the article's own macro definitions (\texttt{\textbackslash E}, \texttt{\textbackslash Prob}, \texttt{\textbackslash TV}) and the theorem environments the retained lines use; the article ships no class or style file, so none is shipped or input; the article's own \texttt{thmtools}/\texttt{thm-restate} packages are loaded so that the retained restatable statement of its parameterised main theorem, which the counted statement cites, can be delivered unchanged; sequential numbering of the retained environments in order of appearance, whereas the article prints them as Proposition E.1 and the supporting lemmas of its sections 4--6; equation numbering is not reproduced, and every cross-reference inside the retained text resolves to the wrapper's numbering. No citation occurs inside any retained line, so the wrapper carries no bibliography. No asset-inclusion command, no drawing environment and no image asset occurs in this package.


\section*{Retained context: thm:main\_parameterized (source0.tex lines 655--701)}

\begin{restatable}{theorem}{MainParameterized}
  \label{thm:main_parameterized}

For every $p>0$, $\rho\in(0,1/4)$, and constant $a\ge1$, there exists
  $c=c(p,\rho,a)>0$ such that the following holds.
  For every $R>0$ and $\gamma\in(0,R/2)$, there exists
  $d_0=d_0(p,\rho,a,R/\gamma)$ such that, for every integer
  $d\ge d_0$ and every $\varepsilon_{\rm err}\in(0,1]$, there exists a distribution
  $\mathcal D$ over pairs
  $(\pi,\{\widehat s_\sigma\}_{\sigma>0})$ on $\mathbb R^d$
  satisfying:
  \begin{enumerate}[label=(\arabic*)]
    \item $\pi=\pi_{\rm pre}*\mathcal N(0,\gamma^2I_d)$ for some
      $\pi_{\rm pre}$ supported on $[-R,R]^d$.

    \item For every $\sigma>0$,
      \[
        \E_{X\sim \pi_\sigma}
        \big[\|\widehat s_\sigma(X)-s_{\pi,\sigma}(X)\|_2^p\big]
        \le \frac{\varepsilon_{\rm err}^p}{\sigma^p}.
      \]

    \item For every $\sigma>0$, both $s_{\pi,\sigma}$ and
      $\widehat s_\sigma$ are globally $L_\sigma$-Lipschitz with
      \[
        L_\sigma\le \frac{1}{\gamma^2+\sigma^2}
        +\frac{5R^2d}{(\gamma^2+\sigma^2)^2}.
      \]

    \item Writing
      \[
        H:=\log(d/\varepsilon_{\rm err})+\log(R/\gamma),
        \qquad
        B_d:=\frac{d\log(R/\gamma)}{\sqrt{dH}+H},
      \]
      every adaptive algorithm that, on every execution path, makes at
      most $d^a$ score queries in total and uses at most $cB_d$ distinct
      smoothing levels (distinct values of $\sigma$) satisfies
      \[
        \Pr_{(\pi,\{\widehat s_\sigma\})\sim\mathcal D}
        \Big[\TV\big(\widehat X,\pi\big)\ge 1-\rho\Big]
        \ge 1-\rho.
      \]
      In particular, the same conclusion holds for every adaptive
      algorithm making at most $cB_d$ score queries in total.
  \end{enumerate}
    \end{restatable}


\section*{Retained context: lem:dense-score (source0.tex lines 1377--1386)}

\begin{lemma}[Dense codebooks are score-close]\label{lem:dense-score}
There are universal constants \(c,C>0\) such that, for every
\(\tau>0\) and every integer \(n\ge1\) with
\(\Gamma:=\log n-I_\tau>0\),
\[
   \E_S\mathcal M_{2,\tau}(S)
   \le C G_\tau^2
      \exp\left[-c\min\left\{\frac{\Gamma^2}{d},\Gamma\right\}\right].
\]
\end{lemma}


\section*{Retained context: prop:blind-output (source0.tex lines 1538--1550)}

\begin{proposition}[Level-sensitive null coupling]
\label{prop:blind-output}
Let \(\mathcal A\) be any randomized adaptive algorithm making at most
\(Q\) queries and using at most \(\ell\) distinct noise levels on every execution
path.  Couple its planted and null executions using the same internal
randomness \(\Xi\), independent of \(S\) and including its output
randomness, and let \(\mathsf{Div}\) be the event that the two executions
differ.  Then
\[
   \Prob_{S,\Xi}(\mathsf{Div})
   \le\frac{C\ell w}{m_d}+Q\delta_d.
\]
\end{proposition}


\section{The counted unit: Proposition E.1 and its complete original proof}

\begin{proposition}[Robustness to polynomially bounded reweighting]
\label{prop:bounded-reweighting}
Fix a constant \(C_{\mathrm{wt}}\ge0\).  For each
\(1\le n\le n_{\max}\), fix a deterministic probability vector
\(\eta^{(n)}=(\eta^{(n)}_1,\ldots,\eta^{(n)}_n)\) satisfying
\[
   n\|\eta^{(n)}\|_\infty\le d^{C_{\mathrm{wt}}}.
\]
Replace \(\nu_S\) in the hard construction by
\(\nu^\eta_S:=\sum_{i=1}^n\eta^{(n)}_i\delta_{Y_i}\), and let
\(\mathcal D^\eta\) be the induced weighted hard law.  Then
\(\mathcal D^\eta\) witnesses Theorem~\ref{thm:main_parameterized}
with the same \(B_d\); the constants \(c\) and \(d_0\) may additionally
depend on \(C_{\mathrm{wt}}\).  The result also holds for a random
weight family drawn independently of \(U\) and all codeword locations
and satisfying the
displayed bound almost surely.
\end{proposition}

\begin{proof}
As before, it suffices to consider \(p\ge2\).
Fix \(n\), abbreviate \(\eta:=\eta^{(n)}\), and put
\(\eta_{\max}:=\max_i\eta_i\).
Use the weighted mixture in the likelihood ratio and oracle, enlarging
the constant in \(w\) to depend also on \(C_{\mathrm{wt}}\).
For every fixed \(x\),
\[
   e^{r^\eta_{S,\tau}(x)}
   =\sum_i\eta_i e^{\imath_\tau(Y_i,x)},
   \qquad \E_S e^{r^\eta_{S,\tau}(x)}=1.
\]
The sparse-side bound follows by Markov's inequality.

For mask coverage, generate a target sample as \(X=Y_J+\tau Z\),
where \(J\sim\eta\).  The self term gives
\(r^\eta_{S,\tau}(X)\ge\imath_\tau(Y_J,X)+\log\eta_J\), and
\(\Prob[\log(n\eta_J)<-t]\le e^{-t}\): the small weights have
total mass at most \(n(e^{-t}/n)\).
Combining this with information concentration, taking both deviations
to be \(h/2\), gives, for every fixed \(S\) and \(h\ge0\),
\[
   \pi^\eta_{S,\tau}
   \bigl(r^\eta_{S,\tau}+\log n<I_\tau-h\bigr)
   \le C e^{-c\min\{h^2/d,h\}}.
\]
This bound gives mask coverage and the target witness estimate
\(\E_{\pi^\eta_S}g^\eta_S\ge1-\rho/2\), with the same \(m_d\)
and \(g^\eta_S:=\min\{1,(r^\eta_{S,\gamma})_+/m_d\}\).
The mean-one likelihood identity gives \(\E_Sg^\eta_S(x)\le1/m_d\).

On the dense side, the variance bound \(T/n\) in the proof of
Lemma~\ref{lem:dense-score} becomes
\[
   T\sum_i\eta_i^2\le T\eta_{\max}.
\]
The same truncation argument therefore applies with information gap
\[
   \log(1/\eta_{\max})-I_\tau
   \ge\log n-I_\tau-C_{\mathrm{wt}}\log d.
\]
Since \(H\ge\log d\), enlarging the constant in \(w\) absorbs this
loss and gives the dense-side bound with failure probability
\(\delta_d\).  Proposition~\ref{prop:blind-output} and the witness
argument now prove the theorem with a smaller resource constant \(c\).
All bounds are uniform over the allowed weight families, so conditioning
on an independent random weight family and then averaging proves the
last assertion.
\end{proof}

\end{document}
